\section{引言：课程概述}

本讲是 MIT 6.S191（深度学习导论）2025 年春季的第七讲，由来自 Google 的客座讲师主讲。讲师在 Google 领导 Gemini 应用研究团队（Gemini Applied Research Group），负责将 Gemini 模型部署到 Google 各产品中，同时也在 UC Berkeley 数据科学硕士项目教授机器学习课程。本讲内容源自其在 Google 内部的 LLM Boot Camp 培训课程。

本讲覆盖四大主题：
\begin{enumerate}
    \item \textbf{LLM 101}：语言模型的基本原理与工作机制
    \item \textbf{LLM 的实际应用}：如何从语言模型构建聊天机器人
    \item \textbf{Prompt Engineering}：提示工程的核心技巧
    \item \textbf{Beyond Prompts}：超越提示词的技术——微调、对齐与 AI Agent
\end{enumerate}

讲师的核心教学目标是让学生理解 LLM 的工作方式，建立关于何时以及如何使用 LLM 的直觉，并了解常见的陷阱（即文献中所称的 ``jagged frontier''——参差不齐的能力边界）。

